\section{Solving the graph equation on a cover}
\label{sec:graph-existence}

The scalar identity alone does not construct a graph. We now solve a
nonlinear equation prescribing its divergence in terms of the error from
a given function. The strict tangent band controls that error, while
uniqueness will allow the solutions to descend through the flat-bundle
transition maps and remain invariant under previously added circles.

The graph regularization and gradient-estimate strategy was inspired by
the Jang-equation capillarity method of Schoen--Yau
\cite[Section~4, pp.~251--253]{SchoenYau1981}.  The variable flux, tangent
forcing, stabilization identity, and fixed-rank shortening estimates
needed here are developed in this paper; their existence theorem is not
being applied to the present equation.

We specify the uniformity needed for the nonlinear equation. Let $B$ be a Riemannian cover of a closed manifold, with lifted reference metric $g_{\rm ref}$. Fix nested coordinate charts obtained by lifting a finite collection of relatively compact reference charts, with the smaller charts covering $B$. For $0<\alpha<1$, uniform $C^{j,\alpha}$ norms are the suprema of the usual chart norms on these smaller charts; the larger charts provide room for interior estimates. Equivalent choices give equivalent norms. A metric $g$ is \emph{uniformly controlled} if it is uniformly equivalent to $g_{\rm ref}$ and its coordinate coefficients have bounded derivatives of every order. The metric need not be invariant under deck transformations. We write $C_b^{j,\alpha}$ for these uniform H\"older spaces and $C_b^\infty$ for smooth functions with bounded derivatives of every order.

Products of such covers with any fixed finite torus have the same properties. In all assertions in this Section, constants may depend on the particular finite-dimensional geometry and the fixed equation coefficients.

\begin{lemma}[A uniform linear inverse]\label{lem:uniform-linear}
Suppose that, in the reference charts, a scalar operator has the form
\[
 L\xi=-a^{ij}\partial_i\partial_j\xi+b^i\partial_i\xi+c\xi,
\]
with uniformly $C^{0,\alpha}$ coefficients, uniformly bounded drift, uniformly positive definite principal matrix, and $c\ge c_0>0$. Then
\[
 L:C_b^{2,\alpha}(B)\longrightarrow C_b^{0,\alpha}(B)
\]
is an isomorphism. Its inverse is bounded in terms of the indicated uniform bounds. If the coefficients have bounded derivatives of every order, the inverse also maps $C_b^{j,\alpha}$ boundedly into $C_b^{j+2,\alpha}$ for every $j\ge0$.
\end{lemma}

\begin{proof}
It suffices to prove the assertion when $\dim B\ge2$. If $d=\dim B<2$,
put $T_0=\T^{2-d}$ with its flat product metric and consider
$\overline L=L-\Delta_{T_0}$ on $B\times T_0$, with right-hand side
pulled back from $B$. Its principal matrix is uniformly positive and its
zeroth-order coefficient is still at least $c_0$. The result in dimension
two gives a unique solution; translation invariance of the data and
uniqueness make that solution independent of $T_0$. It therefore descends
to the required solution on $B$. Restriction to a fixed torus point gives
the same uniform estimates, and lifting two possible solutions proves
uniqueness on $B$.

Take a smooth compact exhaustion and solve the zero-boundary Dirichlet problem on each member. Bounded-domain elliptic Dirichlet theory applies since the principal matrix is uniformly positive and the zeroth-order term has the indicated sign; see, for example, the linear Schauder theory in \cite[Theorems~6.2 and~6.14]{GilbargTrudinger2001}. The maximum principle gives, for right-hand side $F$,
\begin{equation}\label{eq:uniform-inverse-sup}
 \norm{\xi}_{\infty}\le c_0^{-1}\norm{F}_{\infty},
\end{equation}
independently of the exhaustion domain. On every smaller reference chart whose enlargement is contained in that domain, the interior estimate gives
\[
 \norm{\xi}_{C^{2,\alpha}(\text{smaller chart})}
 \le C\bigl(\norm{F}_{C^{0,\alpha}(\text{larger chart})}
                         +\norm{\xi}_{\infty}\bigr),
\]
with one constant for all charts. Local compactness and a diagonal subsequence produce a solution on $B$. Each fixed enlarged chart eventually lies in the exhaustion, so the same estimate holds at every point of the limit. Taking the supremum over charts proves that the solution belongs to $C_b^{2,\alpha}$ with a bounded norm. No estimate at the moving exhaustion boundary is needed. If $B$ is compact, one instead uses the closed-manifold version of the same theory.

For uniqueness, suppose $L\xi=0$ and $S=\sup_B\xi>0$. Center uniform charts at a sequence of points where $\xi\to S$. The pulled-back coefficients have a convergent local $C^{0,\beta}$ subsequence, and the functions a convergent $C^{2,\beta}$ subsequence, for $0<\beta<\alpha$. At the center of the limiting chart the limiting function attains the maximum $S$. There its gradient vanishes and its Hessian is negative semidefinite, so
\[
 L\xi\ge c_0S>0,
\]
a contradiction. Apply the same argument to $-\xi$. This also proves the global maximum principle used in \eqref{eq:uniform-inverse-sup} for any solution in the uniform class. Higher-order interior estimates give the final assertion.
\end{proof}

\begin{theorem}[Uniform graph solvability]\label{thm:uniform-graph}
Let $(B,g)$ be uniformly controlled as above, and let $\ell\in C_b^\infty(B)$. Fix $\eps,w,B_0,s>0$ and $0\le\beta<1$. Suppose $N$ is smooth, positive, nondecreasing, constant near zero, and equal to $1$ for $t\ge1/2$. Define the quantities in \eqref{eq:graph-quantities} and
\[
 \lambda(r)=B_0(r^2+s^2)^{\beta/2}.
\]
There is a unique solution in the uniform $C^{2,\alpha}$ class with a uniform band margin to
\begin{equation}\label{eq:graph-equation}
 -\diver_g U
   =\lambda(r)\tan\left(\frac{\ell-\eps f}{w}\right),
 \qquad \abs{\ell-\eps f}<\frac{\pi w}{2}.
\end{equation}
This solution lies in $C_b^\infty(B)$, and the band has a positive uniform margin.
The metric \eqref{eq:graph-metric-update} is uniformly controlled and complete.

The solution is invariant under every isometry preserving the equation's
data. The following parameter statement will be used below. Fix the
reference cover, its metric and nested charts, and the equation parameters
$\eps,w,B_0,s,\beta,N$. Let $(g_a,\ell_a)$ be a jointly smooth family over
a finite-dimensional parameter manifold $A$. For every relatively compact
subset $V'\Subset V$ of a parameter chart, assume common constants
$c_{V'},C_{V'}>0$ such that
\[
 c_{V'}g_{\rm ref}\le g_a\le C_{V'}g_{\rm ref}\qquad(a\in V'),
\]
and common bounds in the fixed reference charts for every mixed derivative
of the coefficients of $g_a$ and of $\ell_a$, including order zero. Then
$a\mapsto f_a$ is smooth in every $C_b^{j,\alpha}$ norm, for $j\ge0$ and
$0<\alpha<1$. If $A$ is compact,
all mixed derivatives of $f_a$ are uniformly bounded and the tangent band
has a common positive margin. For a bundle of such covers, the local
solutions glue whenever fixed reference deck transformations carry one
local data set to another.
\end{theorem}

\begin{proof}
We may assume $\dim B\ge2$. Once the theorem is proved in those dimensions,
the lower-dimensional cases follow by taking the product with
$T_0=\T^{2-\dim B}$, using the product metric and pulling back $\ell$.
The solution on $B\times T_0$ is invariant under $T_0$ translations by
uniqueness, and hence is pulled back from a function on $B$. For such a
function the flux is $(U,0)$ and its product divergence is $\diver_g U$,
so the descended function solves exactly \eqref{eq:graph-equation}.
Restriction to a fixed torus point gives the uniform and parameter bounds,
and lifting two possible solutions gives uniqueness. The graph metric
claims then follow from the same final bounds below.

Replace $\ell$ by $\tau\ell$, $0\le\tau\le1$, and start at $f=0$ for $\tau=0$. We use $C_b^{2,\alpha}(B)$, for a fixed $0<\alpha<1$, as the solution space.

\smallskip\noindent\textbf{Linearization and openness.}
As a function of the covector $df$, the flux $U$ has transverse eigenvalue $r(t)/t$ and radial eigenvalue
\begin{equation}\label{eq:graph-flux-ellipticity}
 r'(t)=\frac{N(t)(1+\mathfrak p(t))}
                {(1+N(t)^2t^2)^{3/2}}>0.
\end{equation}
Both extend positively and smoothly through $t=0$, since $N$ is constant there. Hence the equation is uniformly elliptic on every bounded gradient range.

On the open set of $C_b^{2,\alpha}$ functions with a positive uniform margin in the tangent band, the equation is a smooth map into $C_b^{0,\alpha}$. The apparent nonsmoothness of $r=\abs{U}$ at zero causes no problem: $\lambda$ is a smooth function of $r^2$, and the flux is smooth in $df$. At an individual solution, moving the right-hand side to the left gives the linearization
\begin{equation}\label{eq:graph-linearization}
 -\diver_g\bigl(a_f(d\xi)\bigr)+\mathbf b(d\xi)
   +\frac{\eps\lambda(r)}w
       \sec^2\left(\frac{\tau\ell-\eps f}{w}\right)\xi.
\end{equation}
Here $a_f$ is the derivative of the flux. After expansion, the coefficients are uniformly H\"older, and the zeroth-order coefficient is at least
$\eps B_0s^\beta/w>0$. The individual $C_b^{2,\alpha}$ bound supplies a bounded gradient range, so \cref{lem:uniform-linear} makes the linearization invertible. The Banach implicit function theorem proves openness in $\tau$.

\smallskip\noindent\textbf{Comparison and individual regularity.}
Two solutions in this class agree. To see this, take a positive supremum of their difference, if one exists, and pass to a limiting uniform chart as in \cref{lem:uniform-linear}. At the limiting maximum their gradients agree and their Hessians are ordered. At that common gradient the right-hand side of the equation is strictly decreasing in height, with derivative at most $-\eps B_0s^\beta/w$. The two equations are therefore incompatible. The same argument compares a solution with a strict upper or lower barrier.

Each individual solution has uniform bounds in every higher derivative before we seek bounds independent of $\tau$. Indeed its given $C_b^{2,\alpha}$ norm and band margin make the equation uniformly elliptic with smooth nonlinearities on a bounded range. In a chart, write its nondivergence form as
\[
 a^{ij}(x,df)f_{ij}=H(x,f,df).
\]
For $z=\partial_k f$, differentiation, initially by difference quotients, gives
\begin{equation}\label{eq:individual-graph-bootstrap}
 \begin{split}
 a^{ij}\partial_i\partial_j z
  +\bigl(a^{ij}_{p_l}f_{ij}-H_{p_l}\bigr)\partial_lz-H_fz
   =H_{x_k}-a^{ij}_{x_k}f_{ij}.
 \end{split}
\end{equation}
The coefficients and source are uniformly $C^{0,\alpha}$ using only the individual $C^{2,\alpha}$ bound. The difference quotients have bounded sup norms, so interior Schauder estimates on smaller charts give $f\in C_b^{3,\alpha}$. Iterating gives $C_b^\infty$. These constants may at this stage depend on the individual solution.

\smallskip\noindent\textbf{Uniform tangent-band barriers.}
For $\delta>0$ consider
\[
 f_\pm=\eps^{-1}\bigl(\tau\ell\pm(\pi/2-\delta)w\bigr).
\]
Their gradient and divergence terms are bounded uniformly in $\tau$ and in small $\delta$. At $f_+$ the tangent term is negative and diverges uniformly to $-\infty$ as $\delta\downarrow0$; at $f_-$ it diverges uniformly to $+\infty$. This uses $\lambda\ge B_0s^\beta>0$. Choose a sufficiently small fixed $\delta$ so that these are strict upper and lower barriers. Comparison gives, for every solution along the homotopy,
\begin{equation}\label{eq:uniform-graph-band}
 f_-\le f\le f_+,
 \qquad
 \abs{\tau\ell-\eps f}\le(\pi/2-\delta)w.
\end{equation}
In particular the heights and the band margins are now controlled independently of $\tau$.

\smallskip\noindent\textbf{A uniform gradient bound.}
Fix one solution along the continuation path. If $\sup_B t\le1$, there is
nothing to prove. Otherwise choose points where $t$ approaches its
supremum. The individual higher bounds just established allow the metric,
data, and solution to converge in uniform charts, in enough derivatives
for the graph equation and its Jacobi identity. The limiting slope attains
its maximum at the center. This passage uses finite bounds for that one
solution; the estimate obtained at the maximum will depend only on the
fixed data and the uniform band margin.

At this maximum, $t>1$, so $N=1$ on a neighborhood. The graph is locally an ordinary product graph, $v=(1+t^2)^{-1/2}$ has a minimum, and $dr=0$. Put
\[
 h_e=\frac1w\left(\lambda+\frac{D^2}{\lambda}\right)
     =\frac{\lambda}{w}
          \sec^2\left(\frac{\tau\ell-\eps f}{w}\right).
\]
Differentiating the right-hand side of the equation along $U$, and using $dr=0$ and $vU(f)=r^2$, the vertical translation identity becomes
\begin{equation}\label{eq:gradient-jacobi}
 \bigl(\Delta_{\rm gr}+\abs{II}^2+\Ric_{g+dz^2}(\nu,\nu)\bigr)v
       =v h_e\tau U(\ell)-\eps h_e r^2.
\end{equation}
The band bound implies fixed constants $0<h_{\min}\le h_e\le h_{\max}<\infty$. The fixed uniformly controlled base has $\Ric_g\ge-C_{\rm Ric}g$. At the minimum of $v$, the left-hand side of \eqref{eq:gradient-jacobi} is at least $-C_{\rm Ric}v$. Since $r^2>1/2$ there and $\abs{U(\ell)}\le\norm{d\ell}_{\infty}$, it follows that
\[
 \frac{\eps h_{\min}}2
 \le v\bigl(C_{\rm Ric}+h_{\max}\norm{d\ell}_{\infty}\bigr).
\]
This bounds $v$ below by a positive constant and hence bounds $t$ above, uniformly in $\tau$. One may add $1$ to the denominator when writing an explicit positive bound. All constants in this final estimate depend only on the fixed data and the already uniform band margin; the individual higher norms used to attain the limiting maximum have disappeared.

\smallskip\noindent\textbf{Uniform higher estimates and closedness.}
After the gradient estimate, the flux is uniformly elliptic on a fixed bounded range. Absorbing the metric volume density, write
\[
 -\partial_i\mathcal A^i(x,\partial f)
       =\mathcal H(x,f,\partial f).
\]
Every coordinate derivative $z=\partial_k f$ satisfies
\begin{equation}\label{eq:graph-divergence-bootstrap}
 -\partial_i\bigl(\mathcal A^i_{p_j}\partial_jz\bigr)
  =\mathcal H_{p_j}\partial_jz+\mathcal H_f z+\mathcal H_{x_k}
       +\partial_i\mathcal A^i_{x_k}.
\end{equation}
The drift coefficient and the nondivergence source are uniformly bounded;
the last term is the divergence of a uniformly bounded vector field.
Interior H\"older regularity for scalar divergence-form equations gives a
uniform $C^{0,\alpha'}$ bound for $z$ on smaller charts, for some
$\alpha'>0$; we use \cite[Theorem~8.24]{GilbargTrudinger2001}. In that
theorem's notation take $b^i=d=0$, retain the bounded drift, and treat
$\mathcal H_fz+\mathcal H_{x_k}$ as scalar data and $\mathcal A^i_{x_k}$ as
divergence data, with the signs determined by
\eqref{eq:graph-divergence-bootstrap}. On a fixed chart these bounded data
satisfy its $L^{r/2}$ and $L^r$ hypotheses for any finite $r>\dim B$.
The estimate uses the local $L^2$ norm of $z$, already controlled by the
gradient bound; it requires no bound for the derivatives of $z$. Thus its
constants depend only on ellipticity, the fixed chart size, the displayed
data bounds, and the gradient bound. We obtain a uniform
$C^{1,\alpha'}$ bound for $f$.

The original equation now has uniformly $C^{0,\alpha'}$ nondivergence coefficients and right-hand side. Interior Schauder estimates give $C^{2,\alpha'}$ bounds, and \eqref{eq:individual-graph-bootstrap} gives higher uniform bounds by iteration. In particular these imply uniform $C^{2,\alpha}$ bounds for the originally chosen $\alpha<1$, by using a bound on one further derivative. Given a sequence of solvable parameters tending to $\tau_*$, local compactness on an exhaustion produces a solution at $\tau_*$ with the same uniform bounds and band margin. The solvable set is therefore closed as well as open, and contains all of $[0,1]$.

\smallskip\noindent\textbf{Symmetries and parameter dependence.}
An isometry preserving the data takes a solution to a solution; uniqueness
gives invariance. More generally, an isometry carrying one set of data to
another carries their solutions to one another by pullback. The uniform
solution class is preserved: covariant H\"older norms for the uniformly
controlled metric are equivalent to the reference-chart norms, and these
intrinsic norms are unchanged under an isometry. In particular all old torus translations preserve a solution when they preserve the input data.

For parameter dependence it is essential to use the stated mixed bounds. If $h(x,a)$ is any input coefficient in a local parameter chart, bounds on its mixed derivatives imply smoothness of $a\mapsto h(\cdot,a)$ in every $C_b^{j,\alpha}$ norm. Explicitly, Taylor's formula in $a$, applied also to $x$ derivatives through order $j+1$, gives
\[
 \norm{h(\cdot,a+u)-h(\cdot,a)-D_a h(\cdot,a)u}_{C_b^{j,\alpha}}
       \le C\abs{u}^2.
\]
The additional $x$ derivative controls the H\"older seminorm on the uniformly bounded charts. Applying the same argument to every parameter derivative proves the claimed Banach-space smoothness.

The equation is therefore a smooth Banach map in $a$ as well as $f$.
The fixed reference charts supply the Banach spaces, while the common
metric equivalence on $V'$ supplies locally uniform ellipticity on bounded
gradient ranges. At each parameter,
the invertible linearization gives a locally smooth solution by the
implicit function theorem. Applying the higher-order part of
\cref{lem:uniform-linear} gives the same conclusion in every higher norm,
and uniqueness identifies all the local branches. Compactness of $A$
supplies uniform mixed bounds and a common band margin using finitely many
such neighborhoods.

For gluing, suppose a fixed reference deck transformation $\psi$ carries
one local data set to another, so that
$g_a^{(2)}=\psi^*g_a^{(1)}$ and $\ell_a^{(2)}=\psi^*\ell_a^{(1)}$.
It is then an isometry between those data sets. Naturality and uniqueness
give $f_a^{(2)}=\psi^*f_a^{(1)}$. Extending $\psi$ by the identity on old
circle factors gives the compatibility needed for the flat bundles below.

Finally, $df$ is bounded, $N$ is positive and bounded on its attained range, and $v$ is bounded below. Thus $g^+$ is uniformly equivalent to the reference product metric and has bounded coefficient derivatives of every order. Its completeness follows from that equivalence to a complete reference metric.
\end{proof}

We have obtained a smooth graph for each fixed finite set of data, together
with the symmetry and parameter control needed to repeat the construction.
The estimates in the next section will use its equation to bound scalar
loss independently of these potentially large analytic regularity constants.
